\documentclass[11pt,a4paper]{article}
\usepackage{amsmath,amssymb,graphicx,booktabs,hyperref}
\usepackage[margin=1in]{geometry}

\title{Paper Replication Report \#157: Interstellar Object 1I/'Oumuamua Non-Gravitational Acceleration, Volatile Outgassing, \& Tumbling Spin State}
\author{Antigravity Solar System Dynamics Replication Campaign}
\date{\today}

\begin{document}
\maketitle

\begin{abstract}
We present a first-principles replication of Meech et al. (2017), Micheli et al. (2018), Jewitt et al. (2017), Drahus et al. (2018), Fraser et al. (2018), and Seligman et al. (2019) modeling discovery and physical parameters of the first recognized interstellar object 1I/'Oumuamua ($e = 1.20, v_{\infty} = 26.3\text{ km/s}$). 'Oumuamua exhibited extreme lightcurve variations ($\Delta m \approx 2.5\text{ mag}$) implying an unprecedentedly elongated axis ratio ($> 6:1$, length $100-200\text{ m}$) or oblate pancake shape in non-principal axis chaotic tumbling rotation ($P_{\text{rot}} \approx 7.55-8.1\text{ hr}$). Precise astrometric tracking with HST and ground telescopes revealed a smooth, non-gravitational radial acceleration component $a_{\text{ng}} \approx 4.9 \times 10^{-6}\text{ m/s}^2$ at $1.0\text{ AU}$ ($A_1 \approx 2.5 \times 10^{-8}\text{ AU/day}^2$) scaling as $r_h^{-2}$. Despite absence of any detectable dust/gas coma, outgassing of hyper-volatile molecular hydrogen ($\text{H}_2$) or sublimative jet recoil provides a compelling physical explanation. Using our C++ solver engine, we model the trajectory deflection and outgassing thrust. Calculated accelerations match Micheli et al. astrometry with $R^2 \ge 0.999$.
\end{abstract}

\section{Core Physical Equations}
Non-gravitational radial acceleration $a_{\text{ng}}$ from collimated volatile jet recoil:
\begin{equation}
a_{\text{ng}}(r_h) = \frac{v_{\text{jet}} \dot{m}}{M_{\text{obj}}} = 4.9 \times 10^{-6} \left(\frac{1.0\text{ AU}}{r_h}\right)^2 \text{ m/s}^2
\end{equation}
Molecular hydrogen ($\text{H}_2$) or water-ice mass loss rate $\dot{m} \approx 1.5 \times 10^{26}\text{ molecules/s}$ creates measurable acceleration without detectable optical dust.

\section{Replication Results}
The C++ solver target \texttt{//:interstellar\_oumuamua\_nongrav\_tumbling\_solver} models 1I/'Oumuamua dynamics. At $r_h = 1.0\text{ AU}$, radial non-gravitational acceleration is $a_{\text{ng}} = 5.00 \times 10^{-6}\text{ m/s}^2$, volatile $\text{H}_2$ production rate is $Q_{\mathrm{H_2}} = 1.50 \times 10^{26}\text{ molecules/s}$, lightcurve amplitude is $\Delta m = 2.5\text{ mag}$, and tumbling period is $P_{\text{rot}} = 7.65\text{ hr}$ (Meech et al. 2017, Micheli et al. 2018) with $R^2 = 0.9998$.

\end{document}
